\begin{table}[ht]
  \centering \small
  \begin{tabularx}{\textwidth}{@{} l L C C C @{}}
    \toprule
    \textbf{Dominio}   & \textbf{Oggetto di verifica}          & \textbf{Strategia} & \textbf{Priorità} & \textbf{Implementato}  \\
    \midrule
    D0: Discovery      & Inventario endpoint, Shadow \gls{API} & BLACK/GREY         & P0--P1            & \cmark                 \\
    \addlinespace
    D1: Authentication & Autenticazione, lifecycle token       & BLACK/GREY         & P0--P1            & \cmark                 \\
    \addlinespace
    D2: Authorization  & \gls{RBAC}, \gls{BOLA}, \gls{BFLA}    & GREY               & P1--P2            & \cmark                 \\
    \addlinespace
    D3: Data Integrity & Schema, Mass Assignment               & GREY               & P2                & \cmark                 \\
    \addlinespace
    D4: Availability   & Rate limiting, circuit breaking       & GREY/WHITE         & P1--P2            & \cmark                 \\
    \addlinespace
    D5: Observability  & Log, tracing, audit trail             & WHITE              & P3                & \xmark{} (pianificato) \\
    \addlinespace
    D6: Hardening      & Security headers, config. gateway     & WHITE              & P2--P3            & \cmark                 \\
    \addlinespace
    D7: Business Logic & Flussi sensibili, \gls{SSRF}          & GREY               & P2                & \cmark                 \\
    \bottomrule
  \end{tabularx}
  \caption{Tassonomia degli otto domini di assurance: strategia di testing (Box Gradient D3.P3), priorità e stato implementativo nella versione corrente.}
  \label{tab:domini-box-gradient}
\end{table}
